% =====================================================================
%  Chapter 23 review -- the 14-station revision contract.
%  Revision-only: tables, one-liners, pointers back via \secref/\thmref/
%  \defref/\figref/\coderef/\probref.  No re-derivations.
%  Summit: CF 932F Escape Through Leaf (2700) -- a tree rooted at 1; node i
%  carries a_i, b_i; a jump x -> y (y in subtree of x, y != x) costs a_x*b_y;
%  for every node report the cheapest cost to reach ANY leaf.  dp[leaf]=0 and
%  dp[u] = min over descendants y of (a_u*b_y + dp[y]) = a line container
%  queried at x=a_u over lines (slope b_y, intercept dp[y]); the container is
%  built bottom-up and MERGED small-to-large -- exactly 23.2's merge engine.
%  Both official samples hand-replayed: n=3 -> "10 50 0"; n=4 -> "-300 100 0 0".
%  Exports prob:c29-summit, lst:c29-kit, lst:c29-summit.
%  Mixed set is EXACTLY 10, all live-verified, un-owned, sorted ASCENDING by
%  effective difficulty; each hint names its reduction and cites a c29 listing.
%  CF 438D/920F (Ch.19) and 1114F/895E/981E/718C (Ch.20) are OWNED elsewhere --
%  they re-enter ONLY as prose re-treats, never re-linked.
% =====================================================================

\chaptersection{Chapter 23 Summary}

One sentence: \emph{when a range update or a DP transition is \emph{not} a single composable
monoid tag -- a clamp, a modulo, a ``best line at a point'', a fixed-state automaton -- the
segment tree stops summarising and starts \emph{choosing}: it keeps extra bookkeeping and
prunes/settles/descends (\secref{c29-beats}), it indexes the domain by position and keeps a
line per node (\secref{c29-lichao}), it stores a whole \emph{transform} and merges by
composition (\secref{c29-func}), or it rides the DFS with an undo stack (\secref{c29-path}).}
Chapter~19 folded one aggregate over a range; Chapter~20 rewrote a range with a tag;
Chapter~22 tiled a tree path into \(O(\log n)\) of Chapter~19's ranges. Every one of them assumed
the update collapses to \emph{one} value that composes. This chapter is the four engines you reach
for the moment that assumption breaks -- and the boundary that hands the remainder back. Each row
is a pointer, not a proof.

\begin{center}
\renewcommand{\arraystretch}{1.25}
\begin{tabular}{@{}p{0.24\linewidth} p{0.48\linewidth} p{0.20\linewidth}@{}}
\toprule
\textbf{Result} & \textbf{Statement} & \textbf{Where} \\
\midrule
Li Chao descent & A node covering \([l,r]\) keeps the line that wins at the midpoint \(\mathit{mid}\);
the loser can only beat it on \emph{one} side, so it recurses left or right -- one root-to-leaf walk per
insert and per query. & \defref{c29-lichao}, \thmref{c29-lichao-cost} \\
Arbitrary slopes & No hull vertices and no monotone order are stored, so insert/query are \(\bigO{\log C}\)
even when the slopes and query points arrive in any order. & \thmref{c29-lichao-cost}, \figref{c29-lichao} \\
Tree-Li Chao, two rides & \emph{Rollback}: a reversible container on the DFS stack queries a node's
\textsc{ancestor} lines. \emph{Merge}: a per-node Li Chao tree merged small-to-large with a lazy intercept
offset queries a node's \textsc{descendant} lines. & \secref{c29-path}, \defref{c29-path-merge}, \thmref{c29-path-li} \\
Beats trichotomy & A node storing \((\max,\mathrm{smax},\mathrm{cnt},\mathit{sum})\) answers
``how does \(\mathit{chmin}(x)\) touch me?'' three ways: \(x\ge\max\) prune, \(\mathrm{smax}<x\) settle
the top cells in place, else descend. & \defref{c29-beats} \\
Blending bound & The descend-only-when-forced policy amortises to \(\bigO{(n+q)\log n}\)-ish -- the
second-maximum is what bounds the total descents. & \thmref{c29-beats-cost} \\
Chtholly (ODT) & An ordered \code{set} of maximal equal-runs plus a \code{split}; every update is run-wise.
Expected-fast on random data, \emph{not} worst-case. & \defref{c29-odt}, \probref{c29-896c} \\
Composition tree & Store a \emph{transform} (affine map or \((\min,+)\) matrix) per node; merge by
composition (associative, so a segment tree works); a range is the ordered product. & \defref{c29-func}, \thmref{c29-func-mono} \\
Engine map & Li Chao / beats / ODT / composition own the four non-monoid cells; every other cell is
Chapters~19--22 or a forward chapter. & \defref{c29-map}, \secref{c29-boundary} \\
\bottomrule
\end{tabular}
\end{center}

\medskip
{\small\textbf{Algorithms from this chapter.}}
\begin{center}
\renewcommand{\arraystretch}{1.2}
\begin{tabular}{@{}p{0.36\linewidth} l l >{\raggedright\arraybackslash}p{0.20\linewidth}@{}}
\toprule
\textbf{Algorithm} & \textbf{Prep} & \textbf{Per op} & \textbf{Code} \\
\midrule
Li Chao insert / query (min container on \([lo,hi]\)) & \(\bigO{C}\) or lazy & \(\bigO{\log C}\) & \coderef{c29-lichao} \\
CF 631E (Product Sum) -- two Li Chao sweeps & \(\bigO{n}\) & \(\bigO{\log C}\) & \coderef{c29-631e} \\
Rollback Li Chao on the DFS (ancestor lines) & \(\bigO{n\log C}\) & \(\bigO{\log C}\) & \secref{c29-path} \\
Small-to-large Li Chao \emph{merge} + offset (SPOJ COT6) & -- & \(\bigO{\log^2 n}\) amort. & \coderef{c29-cot6} \\
Segment-tree beats \(\mathit{chmin}\) + range sum & \(\bigO{n}\) & amort. \(\bigO{\log n}\) & \coderef{c29-beats} \\
Chtholly tree, four run-wise ops (CF 896C) & \(\bigO{n}\) & expected \(\bigO{\log n}\) & \coderef{c29-896c} \\
\((\min,+)\) matrix compose + segment tree (CF 750E) & \(\bigO{nK^3}\) & \(\bigO{K^3\log n}\) & \coderef{c29-func}, \coderef{c29-750e} \\
Production advanced-array kit (one listing) & -- & -- & \coderef{c29-kit} \\
CF 932F (Summit) -- tree-merged Li Chao DP & \(\bigO{n\log C}\) & \(\bigO{\log C}\) & \coderef{c29-summit} \\
\bottomrule
\end{tabular}
\end{center}

\chaptersection{Master Decision Flowchart}

Post-study ladder: walk top to bottom, stop at the first branch that matches. It decides
\emph{which non-monoid engine} the update or transition needs -- and hands back every cell that
Chapters~19--22 already own, plus two forward chapters: a hull that needs \emph{deletions} or a
fully dynamic tree goes to link-cut (\(\chref{c23}\)); an \emph{offline} batch of range queries
goes to Mo (\(\chref{c21}\)).

\begin{lstlisting}[language={}, basicstyle=\ttfamily\small, numbers=none, frame=single]
  A range update / DP transition that a single composable monoid tag CANNOT express?
  |
  |-- Is the range just "run a fixed-state automaton / linear recurrence"?
  |     +-- YES: function-composition tree (def:c29-func).  Store a TRANSFORM per
  |         node, merge by composition.  Affine map (x -> a*x+b) = a pair; a
  |         best-cost-per-state machine = a (min,+) matrix.  CF 750E lst:c29-750e.
  |         Fold LEFT TO RIGHT -- composition is NOT commutative.
  |         This is Ch.20's lazy tag, generalised (secref:c15-compose).
  |
  |-- Does the update CLAMP (chmin/chmax) or take a MODULO?
  |     +-- YES, adversarial worst-case data:
  |     |      segment-tree beats (def:c29-beats, lst:c29-beats).  Node keeps
  |     |      (max, smax, cnt, sum): PRUNE if x>=max, SETTLE the top cells if
  |     |      smax<x, else DESCEND.  Amortised blend bound thm:c29-beats-cost.
  |     +-- YES but the data / operations are RNG-generated (random):
  |     |      Chtholly tree / ODT (def:c29-odt, lst:c29-896c).  std::set of
  |     |      equal-runs + split; ~40 lines, expected-fast.  NOT worst case
  |     |      (CF 896C was hacked).  Prefer beats if the judge can be adversarial.
  |     +-- the clamp is on a TREE PATH?  Heavy-light (Ch.22, c16) still cuts the
  |            path into chain ranges; wrap a BEATS node in each chain.  See
  |            def:c29-map / secref:c29-boundary -- HLD wraps beats, not replaces it.
  |
  |-- Is the query "best LINE / best value at a POINT x" with ARBITRARY slopes over
  |   a fixed integer domain -- a DP transition  dp[i]=min_j (m_j*i + b_j)?
  |     +-- YES, plain array / 1-D DP:
  |     |      Li Chao line container (def:c29-lichao, lst:c29-lichao).  Insert a
  |     |      line and query the pointwise min in O(log C); no hull vertices, no
  |     |      monotone order needed.  CF 631E lst:c29-631e.
  |     +-- the lines RIDE A TREE?  two mechanisms (secref:c29-path):
  |          |-- query a node over its ANCESTOR lines as the DFS descends?
  |          |    ROLLBACK Li Chao: insert on enter, undo on exit -- a reversible
  |          |    container on the DFS stack.  thm:c29-path-li.
  |          |-- query a node over its DESCENDANT lines, computed bottom-up?
  |               Li Chao tree MERGED small-to-large with a lazy intercept offset
  |               (def:c29-path-merge, lst:c29-cot6).  SPOJ COT6; Summit CF 932F
  |               lst:c29-summit.
  |
  |-- NONE of the above -- you probably do NOT need this chapter:
        a COMPOSABLE monoid range FOLD -> Ch.19 (c13);  a COMPOSABLE range WRITE
        -> Ch.20 (c15);  a STATIC / historical version -> Ch.21 (c19);  a PATH on
        a fixed mutable tree -> Ch.22 (c16);  deletions from the hull or a fully
        DYNAMIC tree -> link-cut Ch.25 (c23);  an OFFLINE batch -> Mo Ch.27 (c21).
\end{lstlisting}

\chaptersection{Common Mistakes Checklist}

\begin{itemize}[nosep, leftmargin=1.4em]
  \item \textbf{Reaching for a monotone CHT deque when the slopes or query \(x\)'s are
        arbitrary.} The Li Chao container is what you use when there is \emph{no} order to
        exploit -- it never stores hull vertices, only the mid-winner. Trying to keep a convex
        hull with a deque and unsorted insertions is the classic WA that only fails on hidden
        tests. \(\defref{c29-lichao}\), \(\thmref{c29-lichao-cost}\).
  \item \textbf{Querying the Li Chao container at an \(x\) outside its domain \([lo,hi]\).} The
        tree is built over a \emph{fixed} integer range; a query point off that range (or a
        negative \(x\) fed to a \([0,C]\) tree without shifting) silently reads the wrong leaf.
        Map negatives into the range before they reach a node, as \(\coderef{c29-cot6}\) does
        (\(\coderef{c29-lichao}\)).
  \item \textbf{Rollback Li Chao that forgets to undo.} The container is \emph{shared} across DFS
        siblings; every line inserted on entry must be popped on exit. A missing undo lets one
        subtree's lines leak into a sibling -- wrong answers only in one branch of the tree.
        \(\secref{c29-path}\), \(\thmref{c29-path-li}\).
  \item \textbf{Merging descendant hulls small-to-large without the lazy intercept offset.} In
        COT6 the cost of a bottom \emph{above} the current node carries a running additive term;
        merging a child tree into the parent's must shift every intercept by \(\Delta\). Skip the
        offset and every merged line is mis-scaled. \(\defref{c29-path-merge}\), \(\coderef{c29-cot6}\).
  \item \textbf{Beats: settling when you must descend.} \(\mathit{apply\_chmin}(x)\) is legal only
        under the precondition \(\mathrm{smax}<x<\max\). When \(x\le\mathrm{smax}\) the second
        maximum is also caught and you must push down and recurse; applying the in-place patch
        there corrupts the sum. \(\defref{c29-beats}\), \(\coderef{c29-beats}\).
  \item \textbf{Beats: not pushing the parent's \(\max\) into the children before descending.} A
        node's recorded \(\max\) can be below what a child still thinks it is; every descent and
        every query must first relax the children to the parent's cap (\(\coderef{c29-beats}\)).
  \item \textbf{Chtholly tree: \code{split} in the wrong order.} Call \(\mathit{split}(r+1)\)
        \emph{before} \(\mathit{split}(l)\); the second split erases and re-inserts a run, which
        invalidates the iterator the first split returned. The \(\mathit{op}=4\) loop then walks a
        dangling range. \(\coderef{c29-896c}\).
  \item \textbf{Trusting the ODT as worst-case.} Its speed is conditional on many clamps keeping
        the run count low; adversarial (clamp-free) data drives it to \(O(n)\) per op. Use beats for
        the robust bound. \(\defref{c29-odt}\), \(\secref{c29-beats}\).
  \item \textbf{Composition trees: folding right-to-left, or with a wrong identity.} Matrix / affine
        composition is \emph{not} commutative -- a range must be folded in left-to-right order, and
        the empty/product identity is the \((\min,+)\) identity (0 on the diagonal, \(\infty\) off
        it). Either slip silently changes \(\mathit{answer}[0][4]\) in CF 750E.
        \(\thmref{c29-func-mono}\), \(\coderef{c29-func}\), \(\coderef{c29-750e}\).
\end{itemize}

\chaptersection{WA Clinic}

Verdicts are symptoms; the causes are the short list this chapter drilled. Diagnose top-down --
every leaf names a listing.

\begin{figure}[H]
\centering
\begin{tikzpicture}[diagram,
    v/.style={boxnode, fill=primaryblue!12, draw=primaryblue, text width=22mm, align=center},
    b/.style={boxnode, fill=difforange!12, draw=difforange, text width=12mm},
    l/.style={boxnode, fill=successgreen!10, draw=successgreen, text width=48mm,
              font=\scriptsize, align=left}]
  \node[v] (verdict) {\textbf{Verdict on the judge}};
  \node[b, below left=8mm and 5mm of verdict] (wa) {\textbf{WA}};
  \node[b, below right=8mm and 5mm of verdict] (mle) {\textbf{TLE / RE}};
  \node[l, below right=1mm and 2mm of wa.south] (wa1)
    {a Li Chao container queried off its domain, or a rollback that leaked a sibling's lines into a
     node's ancestor set (\(\defref{c29-lichao}\), \(\secref{c29-path}\))};
  \node[l, below left=1mm and 2mm of wa.south] (wa2)
    {a composition tree folded out of order, or an affine/matrix merge with the wrong identity --
     the range product is not commutative (\(\thmref{c29-func-mono}\), \(\coderef{c29-750e}\))};
  \node[l, below right=1mm and 2mm of mle.south] (mle1)
    {a beats \(\mathit{chmin}\) applied when \(x\le\mathrm{smax}\) instead of descending -- sum goes
     wrong \emph{and} the amortised bound breaks (\(\defref{c29-beats}\), \(\coderef{c29-beats}\))};
  \node[l, below left=1mm and 2mm of mle.south] (mle2)
    {an ODT on adversarial (clamp-free) data: the run set grows to \(n\), every split walks it --
     the documented 896C hack (\(\defref{c29-odt}\), \(\coderef{c29-896c}\))};
  \draw[flow] (verdict) -- (wa);
  \draw[flow] (verdict) -- (mle);
  \draw[flow] (wa) -- (wa1);
  \draw[flow] (wa) -- (wa2);
  \draw[flow] (mle) -- (mle1);
  \draw[flow] (mle) -- (mle2);
\end{tikzpicture}
\caption*{WA: a Li Chao read outside its domain or a leaky rollback; a non-commutative composition
folded the wrong way. TLE or RE: a beats clamp that should have descended, or an ODT pushed onto
adversarial data -- both turn an amortised/expected bound into a linear walk.}
\end{figure}

\chaptersection{Quick Reference Card}

\begin{quickref}[Chapter 23 at a glance]
\textbf{When one tag cannot summarise the update, the tree stops folding and starts choosing: keep a
line per node, keep max/second-max, keep a whole transform, or ride the DFS with an undo. }

\medskip
\textbf{Li Chao container (23.1).} Over a fixed integer domain \([lo,hi]\), a node stores the line
winning at its midpoint; \(\mathit{insert}\) swaps to keep the mid-winner and recurses the loser to
the one side it can still win; \(\mathit{query}\) walks root-to-leaf taking the best. Insert and
query \(\bigO{\log C}\) with \emph{arbitrary} slopes and query order
(\(\defref{c29-lichao}\), \(\thmref{c29-lichao-cost}\), \(\coderef{c29-lichao}\)).

\medskip
\textbf{Li Chao on a tree (23.2).} Two ways the container rides a tree. \emph{Rollback}: a reversible
container on the DFS stack -- insert on enter, undo on exit -- answers a node over its ancestor lines.
\emph{Merge}: compute bottom-up, merge the children's pointer Li Chao trees small-to-large with a
lazy intercept offset, answer a node over its descendant lines
(\(\defref{c29-path-merge}\), \(\thmref{c29-path-li}\), \(\coderef{c29-cot6}\)).

\medskip
\textbf{Beats + ODT (23.3).} Beats keeps \((\max,\mathrm{smax},\mathrm{cnt},\mathit{sum})\) per node;
\(\mathit{chmin}(x)\) prunes, settles the maximal cells in place, or descends, amortised to
\(\bigO{(n+q)\log n}\)-ish (\(\defref{c29-beats}\), \(\thmref{c29-beats-cost}\), \(\coderef{c29-beats}\)).
The ODT is its random-data twin: a \code{set} of equal-runs, a \code{split}, run-wise ops, short and
expected-fast but not worst-case (\(\defref{c29-odt}\), \(\coderef{c29-896c}\)).

\medskip
\textbf{Composition trees (23.4).} A node stores a \emph{transform} (an affine pair, or a
\((\min,+)\) matrix encoding a fixed-state automaton); the merge is \emph{ordered composition} --
associative, so a segment tree answers a range as the product of its covering nodes, folded
left-to-right (\(\defref{c29-func}\), \(\thmref{c29-func-mono}\), \(\coderef{c29-func}\),
\(\coderef{c29-750e}\)).

\medskip
\textbf{The boundary (23.5).} Li Chao owns best-line-at-a-point; beats owns worst-case clamp/modulo;
ODT owns the same on random data; composition owns fixed-state automata. Composable fold \(\to\)
Ch.~19 (\(\chref{c13}\)); composable write \(\to\) Ch.~20 (\(\chref{c15}\)); static/history \(\to\)
Ch.~21 (\(\chref{c19}\)); path on a fixed mutable tree \(\to\) Ch.~22 (\(\chref{c16}\)); hull with
deletions \(\to\) Ch.~25 (\(\chref{c23}\)); offline \(\to\) Ch.~27 (\(\chref{c21}\))
(\(\defref{c29-map}\), \(\secref{c29-boundary}\)).

\medskip
\textbf{Translation.} ``best \(m_j x + b_j\) at a point, slopes unordered'' \(\to\) Li Chao. ``range
\(\min\)-assign / modulo, worst case'' \(\to\) beats; same on RNG data \(\to\) ODT. ``range = run a
small automaton / recurrence'' \(\to\) composition tree. ``the lines live on a tree'' \(\to\) rollback
(ancestors) or merge (descendants). \emph{Anything composable into one tag} is still Chapters~19--22.
\end{quickref}

\begin{lstlisting}[caption={Chapter 23 kit: the Li Chao container, the beats chmin node, the Chtholly split, and the (min,+) matrix composition, one compact skeleton each, in one place}, label={lst:c29-kit}]
// ===== Li Chao LINE CONTAINER (min) over integers [LO,HI] -- 23.1 =====
struct Line { ll m,b; ll f(ll x) const { return m*x+b; } };
const ll LINF = (1LL<<62);
Line lc[4*MAXN]; bool hasL[4*MAXN];
void linsert(int i,int lo,int hi,Line g){
    int mid=(lo+hi)>>1;
    if(!hasL[i]){ lc[i]=g; hasL[i]=true; return; }
    Line f=lc[i];
    bool left = g.f(lo) < f.f(lo);
    bool middle = g.f(mid) < f.f(mid);
    if(middle) swap(lc[i],g);            // node keeps the mid-winner
    if(lo==hi) return;
    if(left!=middle) linsert(i<<1,lo,mid,g);   // loser may win LEFT
    else             linsert(i<<1|1,mid+1,hi,g); // ... or RIGHT
}
ll lquery(int i,int lo,int hi,ll x){
    if(!hasL[i]) return LINF;
    ll res=lc[i].f(x); if(lo==hi) return res; int mid=(lo+hi)>>1;
    return min(res, x<=mid ? lquery(i<<1,lo,mid,x) : lquery(i<<1|1,mid+1,hi,x));
}

// ===== SEGMENT-TREE BEATS range chmin + range sum -- 23.3 =====
ll bmx[4*MAXN],bsmx[4*MAXN],bcmx[4*MAXN],bsm[4*MAXN];   // max, second max, cnt max, sum
void bpull(int i){ bsm[i]=bsm[i<<1]+bsm[i<<1|1];
    if(bmx[i<<1]>bmx[i<<1|1]){ bmx[i]=bmx[i<<1]; bcmx[i]=bcmx[i<<1]; bsmx[i]=max(bsmx[i<<1],bmx[i<<1|1]); }
    else if(bmx[i<<1]<bmx[i<<1|1]){ bmx[i]=bmx[i<<1|1]; bcmx[i]=bcmx[i<<1|1]; bsmx[i]=max(bmx[i<<1],bsmx[i<<1|1]); }
    else { bmx[i]=bmx[i<<1]; bcmx[i]=bcmx[i<<1]+bcmx[i<<1|1]; bsmx[i]=max(bsmx[i<<1],bsmx[i<<1|1]); } }
void bapply(int i,ll x){ bsm[i]+=(x-bmx[i])*bcmx[i]; bmx[i]=x; }   // smx[i] < x < mx[i]
void bpush(int i){ if(bmx[i<<1]>bmx[i]) bapply(i<<1,bmx[i]);
                   if(bmx[i<<1|1]>bmx[i]) bapply(i<<1|1,bmx[i]); }
void bchmin(int i,int l,int r,int ql,int qr,ll x){
    if(qr<l||r<ql||x>=bmx[i]) return;                    // PRUNE
    if(ql<=l&&r<=qr&&bsmx[i]<x){ bapply(i,x); return; }  // SETTLE
    bpush(i); int m=(l+r)>>1;                            // else DESCEND
    bchmin(i<<1,l,m,ql,qr,x); bchmin(i<<1|1,m+1,r,ql,qr,x); bpull(i); }

// ===== CHTHOLLY TREE (ODT): set of equal-runs + split -- 23.3 =====
struct Rd { ll l,r,v; bool operator<(const Rd&o) const { return l<o.l; } };
set<Rd> odt;
auto split(ll x){                       // ensure a run starts at x; return its iter
    auto it=odt.lower_bound(Rd{x,0,0});
    if(it!=odt.end()&&it->l==x) return it;
    --it; ll l=it->l,r=it->r,v=it->v; odt.erase(it);
    odt.insert(Rd{l,x-1,v}); return odt.insert(Rd{x,r,v}).first;
}                                       // CALL split(r+1) BEFORE split(l)!

// ===== FUNCTION-COMPOSITION MERGE: ordered (min,+) product -- 23.4 =====
// K = automaton states. Identity = 0 on diagonal, INF off it. NOT commutative.
Mat comp(const Mat&X,const Mat&Y){ Mat Z;         // apply X first, then Y
    for(int i=0;i<K;i++) for(int k=0;k<K;k++) if(X.a[i][k]<INF)
        for(int j=0;j<K;j++) if(Y.a[k][j]<INF)
            Z.a[i][j]=min(Z.a[i][j], X.a[i][k]+Y.a[k][j]);
    return Z; }
// affine variant: pair (a,b) meaning x -> a*x+b;  (a1,b1) after (a2,b2) = (a1*a2, a1*b2+b1)
\end{lstlisting}

\chaptersection{Write-From-Memory Drills}

Five signature skeletons, each lifted from a listing already in the chapter. Cover the marked line,
reproduce from memory, then check the answer one line below.

\begin{lstlisting}[caption={Drill 1 -- the Li Chao insert: keep the mid-winner, send the loser to one side (from lst:c29-lichao)}]
int mid=(lo+hi)>>1;
if(!has[i]){ tree[i]=g; has[i]=1; return; }        // empty node takes it
Line f=tree[i];
bool left = g.f(lo) < f.f(lo);
bool middle = g.f(mid) < f.f(mid);
if(middle) swap(tree[i],g);                        // NODE keeps the mid-winner
if(lo==hi) return;
if(left!=middle) insert(i<<1,lo,mid,g);            // loser can only win LEFT
else             insert(i<<1|1,mid+1,hi,g);        // ... or RIGHT
\end{lstlisting}
\noindent{\scriptsize Key: two comparisons per level -- at the left endpoint and the midpoint. The
loser beats the winner on \emph{exactly one} side (the lines cross once), so \(\mathit{left}\ne
\mathit{middle}\) chooses it; that single recursion is the whole \(\bigO{\log C}\)
(\(\defref{c29-lichao}\), \(\thmref{c29-lichao-cost}\)).}

\begin{lstlisting}[caption={Drill 2 -- beats chmin: prune, settle, or descend; push before descending (from lst:c29-beats)}]
if(qr<l || r<ql || x>=mx[i]) return;                 // PRUNE: nothing to clamp
if(ql<=l && r<=qr && smx[i]<x){ sum+=(x-mx)*cnt; mx=x; return; }  // SETTLE top cells
pushdown(i);                                          // else DESCEND
update_chmin(lc,l,mid,ql,qr,x); update_chmin(rc,mid+1,r,ql,qr,x); pushup(i);
\end{lstlisting}
\noindent{\scriptsize Key: the in-place settle \(sum\mathrel{+}= (x-\max)\cdot\mathrm{cnt}\) is legal
only when \(\mathrm{smax}<x<\max\) -- then just the \(\mathrm{cnt}\) maximal cells drop. Push the
parent's cap to the children before you descend, or a child reads a stale maximum
(\(\defref{c29-beats}\), \(\thmref{c29-beats-cost}\)).}

\begin{lstlisting}[caption={Drill 3 -- ODT split: carve a run at a boundary (from lst:c29-896c)}]
auto split(ll x){
    auto it = tr.lower_bound(Node{x,0,0,0});
    if(it!=tr.end() && it->l==x) return it;           // a run already starts here
    --it; ll l=it->l, r=it->r, v=it->val; tr.erase(it);
    tr.insert(Node{l, x-1, v});                       // left piece keeps the old value
    return tr.insert(Node{x, r, v}).first;            // right piece starts at x
}
\end{lstlisting}
\noindent{\scriptsize Key: erase the one run that \emph{straddles} \(x\) and re-insert its two
halves. In the driver, always \code{split(r+1)} \emph{before} \code{split(l)} so the returned
iterator for \(l\) is not invalidated (\(\defref{c29-odt}\), \(\coderef{c29-896c}\)).}

\begin{lstlisting}[caption={Drill 4 -- ordered (min,+) composition: the merge of a function tree (from lst:c29-func)}]
Mat compose(const Mat& X, const Mat& Y){   // X THEN Y  (order matters!)
    Mat Z;                                 // all-INF; identity is 0 on diag, INF off
    for(int i=0;i<K;i++) for(int k=0;k<K;k++) if(X.a[i][k]<INF)
        for(int j=0;j<K;j++) if(Y.a[k][j]<INF)
            Z.a[i][j]=min(Z.a[i][j], X.a[i][k]+Y.a[k][j]);
    return Z; }
\end{lstlisting}
\noindent{\scriptsize Key: \((\min,+)\) product reads ``best path \(i\to j\) via some middle
\(k\)''; it is associative (so a segment tree stores it) but \emph{not} commutative, so a range is
folded strictly left-to-right. The identity for an empty range is \(0\) on the diagonal
(\(\defref{c29-func}\), \(\thmref{c29-func-mono}\)).}

\begin{lstlisting}[caption={Drill 5 -- tree-merged Li Chao DP: post-order, merge children, query, insert (from lst:c29-summit / lst:c29-cot6)}]
void dfs(int u){
    if(leaf(u)){ dp[u]=0; }                         // nothing below a leaf
    else { hull[u]=biggest child hull;              // small-to-large base
        for(int v: ch[u]) if(v!=big){ mergeInto(hull[u], hull[v]); }  // fold descendant lines
        dp[u]=qry(hull[u], x=a[u]); }               // best descendant line AT this node's a
    ins(hull[u], m=b[u], c=dp[u]);                  // publish u's own line for its parent
}
\end{lstlisting}
\noindent{\scriptsize Key: merge children \emph{first} (their hulls already hold their subtrees),
read the answer at \(x=a_u\), then insert the line \((b_u,\mathit{dp}[u])\) so the parent sees this
node as one candidate. That is the Summit and COT6 in five lines
(\(\defref{c29-path-merge}\), \(\thmref{c29-path-li}\), \(\probref{c29-summit}\)).}

\chaptersection{Complete Problem Checklist}

Every problem this chapter offers -- four fully solved in-section, worked entries, and one Summit.
Tick \emph{only} after an unaided solve; a hinted attempt counts as half.

\begin{center}
\footnotesize
\setlength{\tabcolsep}{4pt}
\renewcommand{\arraystretch}{1.25}
\begin{tabular}{@{}c p{0.16\linewidth} l p{0.20\linewidth} >{\raggedright\arraybackslash}p{0.26\linewidth} c@{}}
\toprule
\textbf{\#} & \textbf{Problem} & \textbf{Source} & \textbf{Key technique} & \textbf{What to look for} & \textbf{Done?} \\
\midrule
\multicolumn{6}{l}{\textbf{Section 23.1 -- the Li Chao line container}} \\
W23.1 & insert \(2x{+}1,-x{+}6,x{-}2\) on \([1,4]\); read \(x{=}1,3,4\) & worked & \lv{4} mid-winner descent & envelope \((-1,0,1,2)\); loser recurses to one side only & \(\square\) \\
W23.1b & Product Sum & CF 631E & \lv{5} line-in-the-destination DP & move one element; forward+backward Li Chao sweeps & \(\square\) \\
\midrule
\multicolumn{6}{l}{\textbf{Section 23.2 -- Li Chao on a tree}} \\
W23.2 & rollback (ancestors) vs merge (descendants): which fits? & worked & \lv{5} DFS-stack container & enter=insert/exit=undo vs bottom-up small-to-large offset & \(\square\) \\
W23.2b & Cut on a tree & SPOJ COT6 & \lv{6} Li Chao merged up the tree & quadratic chain-tiling; \(\mathit{dp}[1]=42\); lazy intercept offset & \(\square\) \\
\midrule
\multicolumn{6}{l}{\textbf{Section 23.3 -- segment-tree beats and the Chtholly tree}} \\
W23.3 & \(\mathit{chmin}(3)\) on \(A[1..8]\), read sum & worked & \lv{4} prune/settle/descend & \(A{=}(3,2,4,5,1,1,5,3)\): \(24\to19\); descend when \(x\le\mathrm{smax}\) & \(\square\) \\
W23.3b & Willem, Chtholly and Seniorious & CF 896C & \lv{5} ODT: equal-run set + split & LCG repro; op4 power-sum mod; split \(r{+}1\) before \(l\) & \(\square\) \\
\midrule
\multicolumn{6}{l}{\textbf{Section 23.4 -- function-composition trees}} \\
W23.4 & compose two affine pairs / the \((\min,+)\) identity & worked & \lv{4} ordered composition & \((a_1,b_1)\!\circ\!(a_2,b_2)=(a_1a_2,\,a_1b_2+b_1)\); diag 0 & \(\square\) \\
W23.4b & New Year and Old Subsequence & CF 750E & \lv{5} min-plus automaton segtree & 2017 present, 2016 absent; answer \(P[0][4]\); \(-1\) if INF & \(\square\) \\
\midrule
\multicolumn{6}{l}{\textbf{Section 23.5 -- the engine boundary (worked only)}} \\
W23.5 & place four specs on the advanced-array map & worked & \lv{5} which non-monoid engine & best-line/clamp/automaton/tree-riding -- pick the cell & \(\square\) \\
\midrule
\multicolumn{6}{l}{\textbf{Summit (chapter end)}} \\
W23 & Escape Through Leaf & CF 932F & \lv{6} tree-merged Li Chao DP & jump cost \(a_x b_y\); \(\mathit{dp}[u]=\min_y(a_u b_y+\mathit{dp}[y])\); leaves 0 & \(\square\) \\
\bottomrule
\end{tabular}
\end{center}


\chaptersection{Recommended Study Order}

Five focused days. Day~1 internalises the Li Chao container and its one-path descent; day~2 the two
ways the container rides a tree (rollback versus merge); day~3 the beats node and the ODT; day~4 the
function-composition tree and its \((\min,+)\) matrix; day~5 the engine map and the Summit.

\begin{center}
\small
\renewcommand{\arraystretch}{1.25}
\begin{tabular}{@{}llp{0.55\linewidth}@{}}
\toprule
\textbf{Day} & \textbf{Focus} & \textbf{Tasks} \\
\midrule
1 & Read + do 23.1 & hand-insert \(2x{+}1,-x{+}6,x{-}2\) on \([1,4]\), read the envelope \((-1,0,1,2)\) at \(x{=}1,3,4\) (\(\figref{c29-lichao}\)); write \(\coderef{c29-lichao}\) cold; then CF 631E with its two Li Chao sweeps (\(\coderef{c29-631e}\)) \\
2 & Read + do 23.2 & decide rollback-vs-merge for three DP shapes (\(\secref{c29-path}\)); trace COT6 bottom-up to \(\mathit{dp}[1]=42\) on the sample tree (\(\figref{c29-cot6-tree}\)); reproduce the merge-with-offset skeleton (\(\coderef{c29-cot6}\)) \\
3 & Read + do 23.3 & run \(\mathit{chmin}(3)\) on \(A[1..8]\), confirm \(24\to19\), name prune/settle/descend per node (\(\coderef{c29-beats}\)); write the beats \(\mathit{pushup}\) cold; then the ODT of CF 896C and the LCG reproduction (\(\coderef{c29-896c}\)) \\
4 & Read + do 23.4 & compose two affine pairs and the \((\min,+)\) identity by hand (\(\coderef{c29-func}\)); build CF 750E's five-state automaton column by column, fold a range, read \([0][4]\) (\(\coderef{c29-750e}\)) \\
5 & Do 23.5 + Summit & place the four specs on the engine map (\(\defref{c29-map}\)); then the Summit CF 932F -- derive \(\mathit{dp}[u]=\min_y(a_u b_y+\mathit{dp}[y])\) and replay both samples (\(\probref{c29-summit}\)); first 3 mixed \\
\bottomrule
\end{tabular}
\end{center}

\chaptersection{Difficulty Progression}

Tiers are competency gates -- each row names the \emph{skill} the tier unlocks.

\begin{center}
\small
\renewcommand{\arraystretch}{1.3}
\begin{tabular}{@{}lp{0.32\linewidth}p{0.40\linewidth}@{}}
\toprule
\textbf{Tier} & \textbf{Problems} & \textbf{What passing this tier proves} \\
\midrule
\lv{3} \ \(\bigstar\bigstar\bigstar\)
  & the Li Chao insert/query on an \(x\)-grid; the beats \(\mathit{pushup}\) node arithmetic
  & you run the mid-winner descent and combine \((\max,\mathrm{smax},\mathrm{cnt})\) by hand \\
\lv{4} \ \(\bigstar\bigstar\bigstar\bigstar\)
  & the beats \(\mathit{chmin}(3)\) trace, the affine/\((\min,+)\) identity, CF 631E's single sweep
  & you prune/settle/descend correctly and fold an ordered composition with the right identity \\
\lv{5} \ \(\bigstar\bigstar\bigstar\bigstar\bigstar\)
  & CF 896C (ODT), CF 750E (matrix segtree), rollback-vs-merge, mixed at CF~2300--2600
  & you pick the engine, keep a fixed-state automaton as a matrix, and split runs in the right order \\
\lv{6} \ \(\bigstar\bigstar\bigstar\bigstar\bigstar\bigstar\)
  & SPOJ COT6 (Li Chao merge), CF 932F (Summit), the CF~3100 beats/ODT mixed set
  & you merge per-node Li Chao trees bottom-up with an offset and hold two engines in one problem \\
\bottomrule
\end{tabular}
\end{center}

\chaptersection{Reduction Gallery}

Every entry is ``one non-monoid shape \(\to\) one engine.'' Each line collapses a problem family
onto a listing already in the chapter.

\begin{itemize}[nosep, leftmargin=1.4em]
  \item ``best \(m_j x + b_j\) at a point \(x\), slopes and queries in arbitrary order'' \(\to\) the
        Li Chao line container: one insert-descent and one query-walk, both \(\bigO{\log C}\), no hull
        vertices (\(\defref{c29-lichao}\), \(\coderef{c29-lichao}\)).
  \item ``a DP whose transition is a line in the \emph{destination} index'' (CF 631E, move one
        element) \(\to\) sweep left-to-right with a container of source-lines queried at the
        destination \(x\), and again right-to-left (\(\coderef{c29-631e}\), \(\probref{c29-631e}\)).
  \item ``tree DP taking a \(\min\) over \emph{ancestor} lines'' \(\to\) a \textbf{rollback} Li Chao
        container riding the DFS stack -- insert on enter, undo on exit
        (\(\secref{c29-path}\), \(\thmref{c29-path-li}\)).
  \item ``tree DP taking a \(\min\) over \emph{descendant} lines, computed bottom-up'' (SPOJ COT6,
        Summit CF 932F) \(\to\) a per-node Li Chao tree \textbf{merged} small-to-large with a lazy
        intercept offset (\(\defref{c29-path-merge}\), \(\coderef{c29-cot6}\), \(\coderef{c29-summit}\)).
  \item ``range \(\mathit{chmin}/\mathit{chmax}\)/modulo, worst-case data, plus a range sum'' \(\to\)
        beats: keep the maximum, second maximum, count of the maximum and the sum, then
        prune/settle/descend
        (\(\defref{c29-beats}\), \(\thmref{c29-beats-cost}\), \(\coderef{c29-beats}\)).
  \item ``the same clamps but the operations are RNG-generated'' (CF 896C) \(\to\) the Chtholly tree:
        a \code{set} of equal-runs and a \code{split}; expected-fast, not worst-case
        (\(\defref{c29-odt}\), \(\coderef{c29-896c}\)).
  \item ``a clamp on a \emph{tree path}'' \(\to\) heavy-light still supplies the \(O(\log n)\) chain
        ranges; only the leaf engine becomes beats -- HLD \emph{wraps} beats
        (\(\secref{c29-boundary}\), back-citing \(\secref{c16-boundary}\)'s case~(C)).
  \item ``a range = run a fixed-state automaton / linear recurrence'' (CF 750E) \(\to\) store the
        transition as a \((\min,+)\) matrix and fold the range as an ordered product
        (\(\defref{c29-func}\), \(\coderef{c29-func}\), \(\coderef{c29-750e}\)).
\end{itemize}

\chaptersection{Mixed Problem Set}

Ten entries, each live-verified against the current problem sets and absent from every other roster
in the book (the owned IDs -- CF 438D / 920F in Ch.~19, CF 1114F / 895E / 981E / 718C in Ch.~20, and
this chapter's own solved 631E / COT6 / 896C / 750E -- are excluded). None says ``use Li Chao'' or
``use beats'' -- you recognise the engine from \emph{what the update or transition is} and \emph{what
order the data arrives in}. Each hint names its reduction and the Chapter~23 listing or section it
leans on, ordered ascending by effective difficulty.

\begin{mixedset}
  \practice{The Bakery}{CF 833B}
    {https://codeforces.com/problemset/problem/833/B}{CF 2200 L4}
    {A partition DP -- split the array into \(k\) groups maximising the sum of ``number of distinct
     values per group.'' The transition \(\mathit{dp}[g][i]=\max_{j<i}(\mathit{dp}[g{-}1][j]+\mathit{cost})\)
     is the same \emph{shape} as 23.1's best-over-candidates line, but here \(\mathit{cost}\) is not a
     line -- it is a distinct-count you maintain with a moving window. The live tags are \emph{divide
     and conquer}, not geometry: accelerate the DP with divide-and-conquer optimisation (recurse on the
     opt-range) plus a two-pointer BIT for the cost. Recognise \emph{why} it is not a Li Chao problem --
     the candidate value is not affine in the decision variable (\(\secref{c29-lichao}\)'s ``arbitrary
     slope'' vs a data-dependent cost).}
  \practice{Chain Chips}{CF 1814E}
    {https://codeforces.com/problemset/problem/1814/E}{CF 2300 L5}
    {The state you carry along a chain is small and fixed, and each link applies a transition to it.
     That is exactly a function-composition tree (\(\defref{c29-func}\)): store a per-link transform
     and merge by composition -- a tiny matrix / tuple, composed in order, not folded as a scalar.
     The tool from \(\coderef{c29-func}\) carries over verbatim; the only new work is \emph{designing}
     the state so the composition stays closed.}
  \practice{Intersection and Union}{CF 1743F}
    {https://codeforces.com/problemset/problem/1743/F}{CF 2300 L5}
    {Intervals combine, and the quantity you propagate through a merge is a small fixed-dimension
     linear map (a matrix of conditional values). Compose the two intervals' maps with the ordered
     product of \(\coderef{c29-func}\) -- the same associativity that lets a segment tree store a
     \((\min,+)\) automaton in CF 750E (\(\thmref{c29-func-mono}\)). Read the answer off one entry of
     the merged matrix, exactly as \(\coderef{c29-750e}\) reads \([0][4]\).}
  \practice{William The Oblivious}{CF 1609E}
    {https://codeforces.com/problemset/problem/1609/E}{CF 2400 L5}
    {Range queries answered by \emph{composing} per-element transitions, but the state is a bitmask
     over a handful of special letters. Store a small transition table (a map on \(2^{|\text{mask}|}\)
     states) per character and merge by composition -- a function-composition tree over a bit-state
     (\(\defref{c29-func}\)). The \(\coderef{c29-func}\) \code{compose} skeleton is unchanged; only
     the matrix entries differ.}
  \practice{Legen\ldots}{CF 696D}
    {https://codeforces.com/problemset/problem/696/D}{CF 2500 L5}
    {A pattern-matching automaton runs over the string and you need ``how many positions reach an
     accepting state after \(k\) transitions.'' Each character is a fixed-state transition matrix and
     a substring is their ordered product -- \(\coderef{c29-750e}\) verbatim in shape, with the
     automaton of the pattern in place of \(\mathtt{2017}\). Chain the matrices on a segment tree
     (\(\thmref{c29-func-mono}\)); the exponent \(k\) is a repeated composition, so the same
     \(\coderef{c29-func}\) product powers the matrix.}
  \practice{Labyrinth Adventures}{CF 1681E}
    {https://codeforces.com/problemset/problem/1681/E}{CF 2600 L5}
    {Distances through a repeating corridor are a shortest-path DP over a fixed set of vertices -- the
     textbook \((\min,+)\) matrix. Composing two corridor segments is \((\min,+)\) multiplication of
     the type in \(\coderef{c29-func}\) (min over the middle vertex), not ordinary arithmetic; the
     identity is \(0\) on the diagonal and \(\infty\) off it. Fold the whole route as one ordered
     product (\(\defref{c29-func}\), \(\thmref{c29-func-mono}\)).}
  \practice{Little Pony and Lord Tirek}{CF 453E}
    {https://codeforces.com/problemset/problem/453/E}{CF 3100 L6}
    {A range update that drains each value toward a cap and a range query for a ``drained sum.'' A
     single lazy tag cannot summarise a per-element clamp -- that is the wall beats climbs
     (\(\secref{c29-beats}\)). Keep enough per-node bookkeeping that the clamp either prunes, settles
     in place, or descends, and lean on the amortised blend bound \(\thmref{c29-beats-cost}\). The
     \(\coderef{c29-beats}\) \((\max,\mathrm{smax},\mathrm{cnt},\mathit{sum})\) node is the seed; this
     variant adds periodic sampling to the settle step.}
  \practice{Bear and Bad Powers of 42}{CF 679E}
    {https://codeforces.com/problemset/problem/679/E}{CF 3100 L6}
    {Range \(\mathit{chmin}\)-style updates interleaved with a search for the first element hitting a
     special value -- the classic beats discipline (prune/settle/descend, push before you descend,
     \(\defref{c29-beats}\)). The twist is that the predicate is ``is a bad power of 42,'' so augment
     the beats node with a ``nearest special value in range'' field and settle it alongside the max.
     Recognise the engine from the \(\mathit{chmin}\) core, then bolt the query onto it
     (\(\coderef{c29-beats}\)).}
  \practice{Nagini}{CF 855F}
    {https://codeforces.com/problemset/problem/855/F}{CF 3100 L6}
    {Path updates on a tree that \emph{clamp} values toward a colour count, with the tree flattened to
     chains. That is heavy-light (\(\secref{c16-boundary}\)'s path-clamp case) \emph{wrapping} a beats
     node -- the chain ranges come from Ch.~22's decomposition, the non-composable clamp is handled by
     the beats \((\max,\mathrm{smax},\mathrm{cnt})\) engine (\(\coderef{c29-beats}\)). The live tag is
     \emph{binary search}: a query asks for the first position satisfying a predicate, answered by
     walking the beats tree. Two engines, one problem -- exactly the composition \(\secref{c29-boundary}\)
     draws.}
  \practice{Welcome home, Chtholly}{CF 896E}
    {https://codeforces.com/problemset/problem/896/E}{CF 3100 L6}
    {A merge process over runs of equal values -- structurally an ordered-set decomposition like the
     ODT (\(\defref{c29-odt}\)): keep maximal blocks, and let an update collapse many of them into one,
     so the number of blocks you touch amortises small. The same \code{split}-and-walk discipline as
     \(\coderef{c29-896c}\), with a DSU (its live \emph{dsu} tag) to union finished blocks. Reach for
     the ODT only because the process keeps the block count low -- the caution of
     \(\secref{c29-beats}\)'s hacking warning still applies to adversarial input.}
\end{mixedset}


\chaptersection{The Summit -- CF 932F Escape Through Leaf}

\begin{notebox}[Why this is the chapter top rope]
Every in-section engine here reaches for \emph{one} of the chapter's four ideas; the Summit makes you
use \emph{two} at once. The DP is a \(\min\) over lines with \emph{arbitrary} slopes -- Li Chao
(\(\secref{c29-lichao}\)) -- and those lines live on a \emph{tree}, computed bottom-up over each node's
\emph{descendants} -- the merge half of \(\secref{c29-path}\). Read \(\mathit{dp}[u]=\min_{y}
(a_u\,b_y+\mathit{dp}[y])\) and you see a line container queried at \(x=a_u\), whose lines are
\((\text{slope } b_y,\ \text{intercept } \mathit{dp}[y])\) -- but only the lines \emph{below} \(u\). So
you merge the children's Li Chao trees small-to-large with the same machinery as COT6
(\(\defref{c29-path-merge}\), \(\coderef{c29-cot6}\)) and fold this chapter's two signature moves into
a single post-order walk. It is also the honest summit rather than the highest badge: several CF~3100
mixed entries are individually harder variants you solve unaided -- this is the hardest problem the
chapter \emph{fully derives}.
\end{notebox}

\solvedproblem{c29-summit}{Escape Through Leaf}{CF 932F}
  {https://codeforces.com/problemset/problem/932/F}{CF 2700 L5}
  {Li Chao line container merged up a tree (small-to-large); the descendant-min DP folds lines of slope b and intercept dp, queried at the node's a; leaves cost 0}

\begin{problemstatement}[CF 932F Escape Through Leaf]
A tree of \(n\) nodes numbered \(1..n\), rooted at node~\(1\). Node \(i\) carries two integers
\(a_i\) and \(b_i\). From a node \(x\) you may \emph{jump} to any node \(y\) that lies in the subtree
of \(x\) (with \(y\ne x\)); the cost of that single jump is \(a_x\cdot b_y\). A route is a sequence of
one or more jumps, and its cost is the sum of the jump costs.

For \emph{every} node, report the minimum total cost of a route that ends at \emph{any leaf} of the
tree. The root is never considered a leaf even if it has degree one, and you can never jump to
yourself; a node that is already a leaf has cost \(0\).

\textbf{Input.} Line~1: \(n\) (\(2\le n\le 10^{5}\)). Line~2: \(a_1,\dots,a_n\)
(\(-10^{5}\le a_i\le 10^{5}\)). Line~3: \(b_1,\dots,b_n\) (\(-10^{5}\le b_i\le 10^{5}\)). Next
\(n-1\) lines: two endpoints of an edge.

\textbf{Output.} \(n\) integers, node \(i\)'s minimum escape cost.

\textbf{Sample 1.} \(n{=}3\), \(a=(2,10,-1)\), \(b=(7,-7,5)\), edges \((2,3),(2,1)\)
\(\to\) \texttt{10 50 0}.\\
\textbf{Sample 2.} \(n{=}4\), \(a=(5,-10,5,7)\), \(b=(-8,-80,-3,-10)\), edges \((2,1),(2,4),(1,3)\)
\(\to\) \texttt{-300 100 0 0}.
\end{problemstatement}

\begin{attemptgate}[Attempt it yourself first -- \lv{6}, budget a real hour]
Walking every descendant of every node is \(\bigO{n^2}\) (worse with the path lengths) -- \tle at
\(10^5\), and the slopes \(b_y\) are arbitrary, so a monotone convex-hull deque cannot even be built.
\peek{freeze the destination: if \(u\) jumps to a node \(y\) in its subtree and then follows \(y\)'s own best route, what is the total cost in terms of \(\mathit{dp}[y]\)?}
\nudge{That total is \(a_u\cdot b_y + \mathit{dp}[y]\) -- a \emph{line in} \(a_u\) whose slope is
\(b_y\) and intercept is \(\mathit{dp}[y]\). So \(\mathit{dp}[u]=\min_{y\text{ below }u}
(b_y\,x+\mathit{dp}[y])\) evaluated at \(x=a_u\): a Li Chao container queried at one point.}
\fullhint{Compute \(\mathit{dp}\) bottom-up. Each node \(u\) owns a Li Chao tree holding one line
per node of its subtree. Merge the children's trees into \(u\)'s \emph{small-to-large} (the same
pointer merge and offset machinery as COT6, \(\coderef{c29-cot6}\)); the merged tree now holds every
descendant line, so read \(\mathit{dp}[u]=\) query at \(x=a_u\) (a leaf, whose tree is empty, simply
has \(\mathit{dp}=0\)). Then insert \(u\)'s \emph{own} line \((m=b_u,\ c=\mathit{dp}[u])\) so its
parent sees it. Domain of \(x\) is \([-10^5,10^5]\); use 64-bit -- products reach \(10^{10}\) and
path sums \(10^{15}\).}
\end{attemptgate}

\solstep{Understanding}
The route from \(u\) is: pick a landing node \(y\) in the subtree of \(u\), pay \(a_u\cdot b_y\) for
the jump, then follow whatever the cheapest route from \(y\) to a leaf is -- which is exactly
\(\mathit{dp}[y]\). So \(\mathit{dp}[u]=\min_{y}\big(a_u\, b_y + \mathit{dp}[y]\big)\) over all
proper descendants \(y\), and a leaf pays \(0\). Read \(a_u\) as the query coordinate \(x\): every
candidate \(y\) contributes the line \(m x + c\) with \(m=b_y\), \(c=\mathit{dp}[y]\), and
\(\mathit{dp}[u]\) is the lower envelope of \emph{all lines below \(u\)}, sampled at \(x=a_u\). One
line container per node, merged up the tree, is the whole answer.

\solstep{Key Observations}
\begin{observation}
The candidates for \(\mathit{dp}[u]\) are precisely the nodes in the subtree of \(u\) other than
\(u\) itself -- and each already contributes its own \((b_y,\mathit{dp}[y])\) line once \(\mathit{dp}\)
is known bottom-up. So a node's answer is a Li Chao query over \emph{descendant} lines, not ancestor
lines: this is the \textbf{merge} mechanism of \(\secref{c29-path}\), not rollback.
\end{observation}
\begin{observation}
A leaf has no proper descendant, so its minimum over an empty set is taken as \(\mathit{dp}=0\) --
``already at a leaf.'' Every internal node then has a non-empty hull (at least its leaf
descendants' lines) and answers a real query. The root, being non-leaf by the statement, is computed
normally.
\end{observation}
\begin{observation}
The slopes \(b_y\) are in arbitrary order and \(a_u\) can be negative, so no monotone CHT applies;
but Li Chao imposes no order (\(\thmref{c29-lichao-cost}\)). Merging each node's Li Chao tree
small-to-large keeps the total number of line re-insertions at \(O(n\log n)\), so the whole DP is
\(O(n\log^2 n)\) -- a leaf's line is folded upward \(\log n\) times, each fold an \(O(\log C)\) insert.
\end{observation}

\solstep{Think It With Me}
\begin{thinkbox}
\deadend{DFS each node's whole subtree to take the minimum}
{Re-scanning descendants per node is \(\bigO{n^2}\) or worse -- \tle at \(10^5\). The inner minimum is
a \emph{line envelope}, not a plain scan; you must reuse structure across nodes.}
\deadend{A monotone convex-hull trick with a deque}
{The deque only works when slopes and query \(x\)'s arrive sorted. Here the \(b_y\) are arbitrary and
\(a_u\) can be negative, so the ``pointer only moves forward'' invariant breaks -- \textbf{WA} on
unsorted tests. This is precisely the case Li Chao exists for (\(\secref{c29-lichao}\)).}
\deadend{One global Li Chao tree reused for every node}
{A node must see only lines \emph{in its own subtree}, not the whole tree; a shared container leaks
sibling and cousin lines into the query -- \textbf{WA}. The container must be \emph{per node}, merged
up from children (\(\defref{c29-path-merge}\)).}
\aha{Give every node a pointer-based Li Chao tree; in a post-order, \emph{merge} the children's trees
into the largest, \emph{query} at \(x=a_u\) for \(\mathit{dp}[u]\), then \emph{insert} the line
\((b_u,\mathit{dp}[u])\). Leaves query an empty tree, so define \(\mathit{dp}=0\) there. Two chapter
engines -- the line container and the tree merge -- compose into five lines.}
\end{thinkbox}

\solstep{The Approach Ladder}
\begin{center}\renewcommand{\arraystretch}{1.25}
\begin{tabular}{@{}c p{0.24\linewidth} p{0.30\linewidth} l c@{}}
\toprule
\textbf{\#} & \textbf{Approach} & \textbf{Idea} & \textbf{Time} & \textbf{Verdict}\\
\midrule
1 & \lv{2}\ per-node subtree scan & re-walk every descendant for the min & \(\bigO{n^2}\) & \tle \\
2 & \lv{4}\ monotone CHT deque & hull with a moving pointer & \(\bigO{n}\) \emph{if} sorted & \textcolor{warningred}{\ding{55}\ \textbf{WA}} \\
3 & \lv{5}\ one global Li Chao & one container, no subtree filter & \(\bigO{n\log C}\) & \textcolor{warningred}{\ding{55}\ \textbf{WA}} \\
4 & \lv{6}\ \textbf{per-node Li Chao merged small-to-large} & post-order merge, query at \(a_u\), insert \((b_u,\mathit{dp}_u)\) & \(\bigO{n\log^2 n}\) & \ac \\
\bottomrule
\end{tabular}\end{center}
Rung~4 is the Summit: rung~3 is right on the transition but wrong on the \emph{scope} -- a node must
see only its subtree's lines -- and merging per node fixes the scope for one extra logarithm.

\solstep{Solution}
Root the tree at~1 and process nodes in post-order. Each node \(u\) owns the root of a pointer-based
Li Chao tree over \(x\in[-10^5,10^5]\) (split with \(\mathit{mid}=\mathit{lo}+((\mathit{hi}-\mathit{lo})\gg1)\)
so negatives are safe). If \(u\) is a leaf, \(\mathit{dp}[u]=0\) and its tree starts empty. Otherwise
merge every child's Li Chao tree into a base (small-to-large), so the tree holds the lines
\((b_y,\mathit{dp}[y])\) of all proper descendants; read \(\mathit{dp}[u]=\) query at \(x=a_u\). Finally
insert \(u\)'s own line \((m=b_u,\ c=\mathit{dp}[u])\) and hand the tree to the parent. Print
\(\mathit{dp}[1..n]\).

\solstep{Proof of Correctness}
\begin{proofbox}
\textbf{Transition.} Any cheapest route from \(u\) starts with a single jump to some proper
descendant \(y\), costs \(a_u b_y\), and then follows a cheapest route from \(y\), which by definition
costs \(\mathit{dp}[y]\); concatenating is optimal by the prefix-optimal (Bellman) argument, so
\(\mathit{dp}[u]=\min_{y\in\mathit{subtree}(u),\,y\ne u}(a_u b_y+\mathit{dp}[y])\), and a leaf (empty
descendant set) has \(\mathit{dp}=0\). \textbf{Induction.} Process bottom-up: when \(u\) is handled,
each child \(c\)'s tree already contains exactly the lines \((b_y,\mathit{dp}[y])\) for \(y\in
\mathit{subtree}(c)\) (inductive hypothesis, plus \(c\) itself inserted after its query). Merging the
children yields a Li Chao tree holding precisely \(u\)'s proper-descendant lines and nothing else
(merge is set-union of those subtrees, and \(u\)'s own line is inserted only \emph{after} the query),
so \(\mathit{qry}(a_u)\) returns the min over exactly the candidates in the transition. \textbf{Li Chao
soundness} (\(\thmref{c29-lichao-cost}\)): insert/query return the true pointwise minimum of the stored
lines regardless of slope order. Hence \(\mathit{dp}[u]\) equals the transition value, and by induction
every \(\mathit{dp}\) is the true minimum escape cost. \(\Box\)
\end{proofbox}

\solstep{Implementation}
\begin{lstlisting}[caption={CF 932F Summit: a dynamic Li Chao tree per node; post-order, merge the children's trees small-to-large, query the descendant envelope at x=a[u] for dp[u], then insert the line (slope b[u], intercept dp[u]); leaves query an empty tree and get dp=0}, label={lst:c29-summit}]
#include <bits/stdc++.h>
using namespace std;
typedef long long ll;
const ll INF = (4e18);
const int LO = -100000, HI = 100000;      // a[u] lies in this domain
struct Node { ll m,b; Node *l,*r;          // line y=m*x+b, plus two child roots
    Node(ll m_,ll b_):m(m_),b(b_),l(0),r(0){} };
ll eval(Node* t, ll x){ return t ? t->m*x + t->b : INF; }
void ins(Node*& t, int lo, int hi, ll m, ll b){   // Li Chao insert (min)
    if(!t){ t=new Node(m,b); return; }
    int mid = lo + ((hi-lo)>>1);
    bool L = (m*lo+b) < (t->m*lo + t->b);
    bool M = (m*mid+b) < (t->m*mid + t->b);
    if(M){ swap(t->m,m); swap(t->b,b); }          // node keeps the mid-winner
    if(lo==hi) return;
    if(L!=M) ins(t->l, lo,   mid, m, b);           // loser can only win LEFT
    else     ins(t->r, mid+1, hi, m, b);           // ... or RIGHT
}
ll qry(Node* t, int lo, int hi, ll x){             // pointwise min along the path
    if(!t) return INF;
    int mid = lo + ((hi-lo)>>1);
    ll r = t->m*x + t->b;
    return min(r, x<=mid ? qry(t->l,lo,mid,x) : qry(t->r,mid+1,hi,x));
}
Node* mergeTrees(Node* a, Node* b, int lo, int hi){ // fold every line of b into a
    if(!a) return b;  if(!b) return a;
    ins(a, lo, hi, b->m, b->b);                    // b's representative line
    if(lo!=hi){ int mid=lo+((hi-lo)>>1);
        a->l = mergeTrees(a->l, b->l, lo, mid);
        a->r = mergeTrees(a->r, b->r, mid+1, hi); }
    return a;
}
int n; ll a[100005], b[100005], dp[100005]; vector<int> adj[100005], ch[100005];
void rootDfs(int u,int p){ for(int v:adj[u]) if(v!=p){ ch[u].push_back(v); rootDfs(v,u); } }
Node* solve(int u){
    Node* H = 0;
    if(ch[u].empty()){ dp[u]=0; }                   // leaf: no proper descendant
    else {
        for(int v: ch[u]){ Node* Cv = solve(v);      // small-to-large: bigger survives
            if(!H) H = Cv; else H = mergeTrees(H, Cv, LO, HI); }
        dp[u] = qry(H, LO, HI, a[u]);                // best descendant line at x=a[u]
    }
    ins(H, LO, HI, b[u], dp[u]);                     // publish u's line for its parent
    return H;
}
int main(){
    scanf("%d",&n);
    for(int i=1;i<=n;i++) scanf("%lld",&a[i]);
    for(int i=1;i<=n;i++) scanf("%lld",&b[i]);
    for(int i=1;i<n;i++){ int u,v; scanf("%d%d",&u,&v); adj[u].push_back(v); adj[v].push_back(u); }
    rootDfs(1,0);
    solve(1);
    for(int i=1;i<=n;i++) printf("%lld ", dp[i]);   // leaves already 0
    return 0;
}
// Two engines, one post-order: the Li Chao CONTAINER (23.1) supplies ins/qry; the Li Chao MERGE
// up a tree (23.2, def:c29-path-merge) supplies mergeTrees. dp[u] is a query over DESCENDANT
// lines (merge, not rollback). 64-bit everywhere: a*b reaches 1e10, path sums reach 1e15.
\end{lstlisting}

\solstep{Trace on the Official Samples}
\begin{tracebox}[Sample 1: root \(1\), chain \(1\text{--}2\text{--}3\), so \(1\!\to\!2\!\to\!3\)]
Node~3 is a leaf: \(\mathit{dp}[3]=0\), insert line \((m=b_3{=}5,\ c{=}0)\). Node~2's hull after merge
\(=\{5x\}\): \(\mathit{dp}[2]=\mathit{qry}(a_2{=}10)=5\cdot10=50\); insert \((m=b_2{=}-7,\ c{=}50)\).
Node~1's hull \(=\{5x,\ -7x+50\}\): \(\mathit{dp}[1]=\min(5\cdot2,\ -7\cdot2+50)=\min(10,\ 36)=10\).
Output \(\boxed{10}\ \boxed{50}\ \boxed{0}\) -- matching the sample, and the note confirms it: node 2
jumps to 3 for \(a_2 b_3=50\), node 1 jumps straight to 3 for \(a_1 b_3=10\). \(\checkmark\)
\end{tracebox}

\noindent\textbf{Replaying sample 2.} Root~1 with edges \((2,1),(1,3)\) makes \(1\)'s children
\(\{2,3\}\), and edge \((2,4)\) makes \(4\) a child of \(2\); leaves are \(3\) and \(4\), so
\(\mathit{dp}[3]=\mathit{dp}[4]=0\). Node~2's hull \(\{m=b_4{=}-10,\ c=0\}\):
\(\mathit{dp}[2]=(-10)\cdot a_2=(-10)(-10)=100\); insert \((b_2{=}-80,\ 100)\). Node~1's merged hull
holds lines \((b_3{=}-3,0)\), \((b_4{=}-10,0)\), \((b_2{=}-80,100)\);
\(\mathit{dp}[1]=\min(-3\cdot5,\ -10\cdot5,\ -80\cdot5+100)=\min(-15,-50,-300)=-300\). The algorithm
prints \(\boxed{-300}\ \boxed{100}\ \boxed{0}\ \boxed{0}\) -- matching the judge. The \(-300\) is
\emph{load-bearing}: it beats jumping straight to a leaf (\(-50\)) precisely because a negative
slope \(\times\) a positive \(a_1\) plus a cheap sub-route wins -- the min-over-lines sees it, a greedy
``nearest leaf'' would not.

\solstep{Complexity}
\begin{complexitybox}
Rooting is \(\bigO{n}\). Every node's line is inserted once into its own hull and then re-folded
whenever its hull is merged into a strictly larger one -- small-to-large caps that at \(O(\log n)\)
re-insertions per line, each an \(\bigO{\log C}\) Li Chao insert (\(\thmref{c29-lichao-cost}\)) with
\(C=2\cdot10^5\). Total \(\bigO{n\log n\log C}\approx 10^{5}\cdot17\cdot18\approx3\cdot10^{7}\)
operations -- \ac in 3s. Memory is one Li Chao node per line per level it survives:
\(\bigO{n\log C}\) worst case, a few tens of MB under 256 MB; use 64-bit throughout
(\(a\cdot b\) up to \(10^{10}\), \(\mathit{dp}\) up to \(10^{15}\)).
\end{complexitybox}

\begin{hardeningbox}
\begin{itemize}[nosep,leftmargin=1.4em]
  \item \textbf{Query \emph{before} you insert \(u\)'s own line.} \(\mathit{dp}[u]\) is a min over
        \emph{proper} descendants; inserting \((b_u,\mathit{dp}[u])\) first would let \(u\) answer with
        its own line (a self-jump, which the statement forbids) -- \textbf{WA}.
  \item \textbf{Leaves answer \(0\), not \(-\infty\) from an empty query.} Define the empty hull to
        give \(\mathit{dp}=0\) explicitly; otherwise the ``already at a leaf'' case is wrong and
        propagates a bogus large value upward.
  \item \textbf{Split the domain with \(\mathit{mid}=\mathit{lo}+((\mathit{hi}-\mathit{lo})\gg1)\), not
        \((\mathit{lo}+\mathit{hi})/2\).} With a negative \(\mathit{lo}\) the sum can round the wrong
        way; the difference form keeps the recursion well-founded (\(\secref{c29-path}\)'s domain note).
  \item \textbf{Merge is set-union of descendant lines -- and the offset of COT6 is \emph{not} needed
        here.} In 932F a line's intercept is exactly \(\mathit{dp}[y]\), independent of the querying
        node, so a plain merge suffices; reusing COT6's lazy \(\Delta\) would double-count. Contrast
        \(\coderef{c29-cot6}\).
  \item \textbf{64-bit everywhere.} \(a_u b_y\) reaches \(10^{10}\) and a route sums \(O(n)\) such
        terms; \code{int} overflows to a \textbf{WA}, and the \(\mathit{INF}\) sentinel must be a
        \(4\cdot10^{18}\)-scale constant, not \(10^9\).
\end{itemize}
\end{hardeningbox}

\begin{recallbox}[Chapter 23 in ten seconds]
When a range update or DP transition is \emph{not} one composable tag, the tree stops summarising and
starts choosing. ``Best \(m_jx+b_j\) at a point, slopes unordered'' \(\to\) the Li Chao container,
insert and query in \(\bigO{\log C}\) (\(\defref{c29-lichao}\), \(\coderef{c29-lichao}\)); let it ride a
tree by \emph{rollback} for ancestor lines or \emph{small-to-large merge} for descendant lines
(\(\secref{c29-path}\), \(\coderef{c29-cot6}\)). A worst-case clamp/modulo \(\to\) beats with a
\((\max,\mathrm{smax},\mathrm{cnt},\mathit{sum})\) node that prunes/settles/descends
(\(\thmref{c29-beats-cost}\), \(\coderef{c29-beats}\)); the same on RNG data \(\to\) the ODT
(\(\defref{c29-odt}\), \(\coderef{c29-896c}\)). A range that runs a fixed automaton \(\to\) a
composition tree folding \((\min,+)\) matrices in order (\(\defref{c29-func}\), \(\coderef{c29-750e}\)).
Summit CF~932F (\(\probref{c29-summit}\)) merges a Li Chao tree up a tree and reads
\(\mathit{dp}[u]=\min_y(a_u b_y+\mathit{dp}[y])\) at \(x=a_u\) -- two engines, one post-order.
\selftest{
  \item State the Li Chao insert rule: which line does a node keep, and which side does the loser
        recurse to? \hfill\lv{3}
  \item Rollback vs merge -- which one answers a node over its \emph{descendant} lines, and why does
        the other fail here? \hfill\lv{5}
  \item In beats, why must \(\mathit{apply\_chmin}\) only run under \(\mathrm{smax}<x<\max\), and what
        do you do when \(x\le\mathrm{smax}\)? \hfill\lv{4}
  \item Write the \((\min,+)\) composition loop: what is the identity, and why can you not swap the
        two operands? \hfill\lv{4}
  \item In the Summit, why is \(u\)'s own line inserted \emph{after} the query, and why is no COT6
        offset needed? \hfill\lv{6}
}
\end{recallbox}


\chaptersection{Tool Kit Inventory}

The machine assembled through Chapter~23. Earlier chapters' rows are carried forward; this chapter's
are appended. Pointers only.

\begin{center}
\small
\renewcommand{\arraystretch}{1.25}
\begin{tabular}{@{}p{0.23\linewidth} l c p{0.22\linewidth} >{\raggedright\arraybackslash\scriptsize}p{0.18\linewidth}@{}}
\toprule
\textbf{Tool} & \textbf{Since} & \textbf{Lv} & \textbf{Trigger phrase} & \textbf{Where} \\
\midrule
\multicolumn{5}{l}{\itshape Carried forward} \\
Plain segment tree (point update, range fold) & 19 & \lv{3} & ``the build/merge skeleton every engine reuses'' & \chref{c13} \\
Lazy tag as a composition monoid & 20 & \lv{4} & ``the composable case; 23.4 generalises it'' & \chref{c15} \\
Persistent structures & 21 & \lv{5} & ``a static/historical version, not a live non-monoid update'' & \chref{c19} \\
Heavy-light decomposition & 22 & \lv{5} & ``path clamp: HLD wraps a beats node'' & \chref{c16} \\
Link-cut / Euler-tour trees & 25 & \lv{5} & ``the hull needs deletions or the tree moves'' & \chref{c23} \\
Mo's algorithm & 27 & \lv{4} & ``an offline batch of these range queries'' & \chref{c21} \\
\midrule
\multicolumn{5}{l}{\itshape Chapter 23 -- new this chapter} \\
Li Chao line container & 23.1 & \lv{4} & ``best line at a point, slopes unordered'' & \defref{c29-lichao}, \coderef{c29-lichao} \\
DP-with-lines (631E sweep) & 23.1 & \lv{5} & ``the transition is a line in the destination'' & \coderef{c29-631e} \\
Rollback Li Chao on the DFS & 23.2 & \lv{5} & ``query a node's ANCESTOR lines'' & \secref{c29-path} \\
Li Chao merge small-to-large + offset & 23.2 & \lv{6} & ``query a node's DESCENDANT lines'' & \defref{c29-path-merge}, \coderef{c29-cot6} \\
Beats \((\max,\mathrm{smax},\mathrm{cnt})\) node & 23.3 & \lv{4} & ``range chmin/chmax/modulo, worst case'' & \defref{c29-beats}, \coderef{c29-beats} \\
Chtholly tree (ODT) + split & 23.3 & \lv{5} & ``same clamps on RNG data'' & \defref{c29-odt}, \coderef{c29-896c} \\
Ordered \((\min,+)\) / affine composition & 23.4 & \lv{4} & ``a range runs a fixed-state automaton'' & \defref{c29-func}, \coderef{c29-func} \\
The advanced-array engine map & 23.5 & \lv{5} & ``which non-monoid engine, and when to hand back'' & \defref{c29-map} \\
Tree-merged Li Chao DP (Summit) & Summit & \lv{6} & ``min over descendant lines, computed bottom-up'' & \coderef{c29-summit} \\
\bottomrule
\end{tabular}
\end{center}

\chaptersection{Related Topics}

\begin{itemize}[nosep, leftmargin=1.4em]
  \item \chref{c13} (Ch.~19, Segment Trees) --- backward: Ch.~19 owns the plain build / point-update /
        range-fold tree; 23.3/23.4 reuse that exact skeleton, only swapping the leaf (a beats node, a
        transform) -- the tree itself is never re-taught. The modulo-prune CF 438D / 920F stay owned
        there (\(\secref{c13-static}\)) and re-enter only as prose.
  \item \chref{c15} (Ch.~20, Lazy Propagation) --- backward, and a promise \textbf{discharged}: Ch.~20's
        \(\secref{c15-strains}\) handed the non-composable \(\min\)-assign / modulo updates forward with
        ``no beats code there''; \(\secref{c29-beats}\) pays it back. 23.4 is the general case of
        Ch.~20's \(\secref{c15-compose}\) composition monoid -- the transform-at-node \emph{is} a lazy
        tag freed from commutativity worries.
  \item \chref{c19} (Ch.~21, Persistent Structures) --- backward/contrast: persistence serves a
        \emph{static or historical} version; Ch.~21's \(\secref{c19-sparse}\) ``beyond'' row handed the
        live non-monoid updates here. The engines are complements -- a container that mutates in place
        versus a version you never overwrite.
  \item \chref{c16} (Ch.~22, Heavy-Light) --- backward, and a promise \textbf{discharged}: Ch.~22's
        \(\secref{c16-boundary}\) case~(C) (a path \(\mathit{chmin}\)) and its NOTE~5b that dropped
        COT6 as ``a Li Chao problem, forward-cited to Ch.~23'' both land here. HLD supplies the chain
        ranges; \secref{c29-beats} supplies the per-range engine -- the two compose, they do not replace.
  \item \chref{c23} (Ch.~25, Link-Cut / Euler-Tour Trees) --- forward: a Li Chao / hull container that
        must support \emph{deletions} of lines, or a tree whose edges change, is out of scope here; the
        rollback workaround answers ``delete the last thing inserted,'' not an arbitrary line
        (\(\secref{c29-boundary}\)).
  \item \chref{c21} (Ch.~27, Mo's Algorithm) --- forward: a batch of \emph{offline} queries that each
        want a pointwise envelope can be re-ordered and swept instead of maintained online
        (\(\secref{c29-boundary}\)).
  \item \chref{c18} (Ch.~24, Centroid Decomposition) --- forward: \emph{counting} or \emph{optimising
        over all paths} by length/aggregate is centroid's job; here the tree only supplies the subtree
        scope of a line container. \textbf{This is the Next chapter.}
\end{itemize}

\begin{lastword}[Last Word --- Chapter 23]
\emph{``When a single value can no longer summarise what an update does -- a clamp that depends on the
data, a best line at a point with unordered slopes, a range that is really a little automaton -- the
segment tree stops folding and starts choosing: it keeps a line per node and sends the loser one way,
it keeps the maximum and second-maximum and prunes or settles or descends, it keeps a whole transform
and composes, or it rides the DFS with an undo stack and hands a subtree's lines back up.''}

\medskip
Chapter~19 folded one aggregate; Chapter~20 rewrote a range with one tag; Chapter~21 kept history;
Chapter~22 tiled a path into ranges. All four rested on the same crutch -- the update collapses to
\emph{one} value that composes. This chapter was the four times that crutch breaks, and the four
engines that stand in the breach: the Li Chao container (\(\thmref{c29-lichao-cost}\)), which indexes
the domain instead of the array and needs no slope order; segment-tree beats
(\(\thmref{c29-beats-cost}\)), which spends extra node bookkeeping to buy an amortised bound for
\(\mathit{chmin}\) and modulo; the Chtholly tree (\(\defref{c29-odt}\)), beats' short, lucky,
random-data twin; and the function-composition tree (\(\defref{c29-func}\)), which is Chapter~20's tag
grown up into an associative, order-respecting product. The Summit CF~932F (\(\probref{c29-summit}\))
took two of them -- a container and a merge -- and ran them as one post-order walk.

\medskip
Carry three reflexes forward:
\begin{enumerate}[nosep, leftmargin=1.6em]
  \item \emph{If a tag cannot summarise the update, store more state and choose.} Lines \(\to\) Li Chao;
        max/second-max \(\to\) beats; a whole transform \(\to\) composition; a subtree's worth of lines
        \(\to\) rollback or merge (\(\defref{c29-map}\)).
  \item \emph{Order matters where composition replaces addition.} A \((\min,+)\) product folds
        left-to-right (\(\thmref{c29-func-mono}\)); an ODT splits \(r+1\) before \(l\)
        (\(\coderef{c29-896c}\)); a Summit node queries \emph{before} it inserts its own line
        (\(\probref{c29-summit}\)).
  \item \emph{Know where the chapter ends.} Fully dynamic hulls with deletions are link-cut
        (Ch.~25 \(\chref{c23}\)); offline batches are Mo (Ch.~27 \(\chref{c21}\)); counting over all
        paths is centroid (Ch.~24 \(\chref{c18}\)) (\(\secref{c29-boundary}\)).
\end{enumerate}

\medskip
And three habits, drilled until boring: bound your container's domain and shift negatives before they
reach a node; state the amortised or expected bound out loud before trusting beats or the ODT on
adversarial data (896C's hack is the warning); and when you merge a hull up a tree, get the scope
right -- ancestors by rollback, descendants by merge, and the node's own line inserted last. The array
you have mastered since Chapter~19 now has four more personalities; which one you summon is exactly the
question \secref{c29-boundary} taught you to ask.

\medskip
\textbf{Next:} \chref{c18} (Centroid Decomposition) leaves the array-and-update engines of Part~IV
behind and turns to the other half of ``operations over trees'': counting, optimising, and answering
questions that quantify over \emph{all paths} at once -- the divide-and-conquer on the tree's shape
that heavy-light and Li Chao only hinted at.
\end{lastword}
